% !TEX root = ../main.tex
% called by main.tex
\chapter{Análisis de la defectología}\label{ch:defectology}

En esta parte quiero dedicar un espacio a comentar algunos de los defectos de impresión típicos. Con esto me refiero a cualquier anomalía, desde un mínimo artefacto imprevisto en la pieza final, a fallos en la impresión. 

\section{Infra-extrusión o Sobre-extrusión}
\todo{Conseguir e incluir foto de sobre extrusión e infra-extrusión}\\
Este tipo de errores son más complicados de ver. Se pueden detectar en la capa superior de una impresión, donde el relleno se une con los perímetros, y se puede solucionar ajustando el multiplicador de extrusión (\texttt{Extrusion Multiplier}) en PrusaSlicer, en la dirección correcta. Si se ve sobre-extrusión, habría que reducirlo ligeramente, y si se ve infra-extrusión, habrá que aumentarlo ligeramente. Es importante comentar que este parámetro vale 1 por defecto, y que un filamento bien calibrado tiene el multiplicador de extrusión oscilando entre 0.94 y 1.04. Valores mayores o menores son indicativos de que el problema es otro.  

\section{Perímetros externos primero}

\section{Pie de elefante (Elephant's foot)}

\section{La pieza se levanta de la cama}

\section{Mal alineamiento de las capas}
